% Lab 3 -- Exercise 1  |  Nguyễn Quốc Đạt – 2510287  |  Original: Ch.3 Ex.3.1 + Ex.3.2 (merged)  |  Status: Adjust
\exercise{Saturation Mode and DC Sweep}{Adjust}{Ch.3 Ex.3.1 + Ex.3.2 (merged)}
\statusreason{\todo{Explain what was changed (component, model, value or merging) and why the learning objective is preserved.}}

\task{\todo{Problem statement: what is given and what must be found (see Ch.3 Ex.3.1 + Ex.3.2 (merged)).}}

\subsubsection*{Circuit}
\circuitfig{ex01}{Circuit of Exercise 3.1}
\components{\todo{List every component value and model, e.g. R1 = 1\,k$\Omega$, D1: 1N4148.}}

\subsubsection*{Theoretical Calculation}
\answerspace{8}
\begin{solution}
\todo{Worked solution with formulas and steps.}
\end{solution}

\subsubsection*{Simulation}
\todo{Describe the LTspice setup: sources, analysis, what to plot or measure.}
\begin{spice}
* TODO: replace with the directives of this exercise
.op
\end{spice}
\simfig{ex01_sim}{Simulation result of Exercise 3.1}

\subsubsection*{Theory vs. Simulation}
\begin{comparison}{Theory vs. simulation, Exercise 3.1}
\cmprow{$I_B$}{mA}{}{}{}
\cmprow{$I_C$}{mA}{}{}{}
\cmprow{$V_{CE}$}{V}{}{}{}
\end{comparison}
\answerspace{3}
\begin{solution}
\todo{Explain the deviations (device model, assumptions such as $V_D = 0.7$\,V, $\beta$, ...).}
\end{solution}
